% Minimal preamble for P2 cards taken from sources/comparability-partitions-2408.13523v2/source0.tex
% (arXiv:2408.13523v2, "Partitioning perfect graphs into comparability graphs").
% Every line below is copied mechanically from that file; nothing was invented, renamed or redefined.
%   line 2    : \documentclass[12pt]{article}
%   line 3    : \usepackage{amssymb,latexsym}
%   line 4    : \usepackage{amsmath}
%   line 5    : \usepackage[dvipdfm]{graphicx}  -- the dvipdfm driver option is dropped, because the
%               assembler forces the xelatex engine and a forced driver option contradicts it.
%   line 16   : \usepackage{amsfonts}
%   line 17   : \usepackage{amsthm}
%   lines 20-30: the author's theorem declarations
%   lines 32-41: the author's own macro definitions
% Omitted source preamble lines and why they are not needed for these body spans:
%   line 1 (provenance comment), lines 6-11 (tikz and its libraries: no drawing in the retained spans),
%   lines 7 and 9 (verbatim), lines 12-13 (page geometry), line 14 (rawfonts),
%   line 15 (psfrag), line 18 (hyperref with the pagebackref option: the retained spans print no
%   citation anchors and the assembly adds its own hyperref machinery), lines 43-61 (the author's
%   backref neutralisation and his margin-note macros, which depend on the omitted hyperref).
% The retained spans of the body use only \em, \bf, \cal, amsthm environments, amsmath displays and
% the macros defined in lines 32-41 below.
\documentclass[12pt]{article}
\usepackage{amssymb,latexsym}
\usepackage{amsmath}
\usepackage{graphicx}
\usepackage{amsfonts}
\usepackage{amsthm}

\newtheorem{theorem}{Theorem}%[section]
\newtheorem{problem}[theorem]{Problem}%[section]
\newtheorem{observation}[theorem]{Observation}
\newtheorem{lemma}[theorem]{Lemma}
\newtheorem{corollary}[theorem]{Corollary}
\newtheorem{definition}{Definition}
\newtheorem{example}[theorem]{Example}
\newtheorem{conjecture}{Conjecture}%[section]
\newtheorem{proposition}[theorem]{Proposition}
\newtheorem{claim}[theorem]{Claim}
\newtheorem{remark}[theorem]{Remark}

\def\qed{\ifhmode\unskip\nobreak\fi\quad\ifmmode\Box\else$\Box$\fi}
\def\pont{\hspace{-6pt}{\bf.\ }}
\def\qed{\hfill $\Box$}

\def\R{{\mbox{\boldmath$\mathcal R$}}}
% the author's second definition of \R, left commented out exactly as in the source
\def\F{{\mathcal F}}
\def\P{{\mathcal P}}

\def\newpage{\vfill\eject}
\title{Every generalized split graph has its edge set partitioned into two comparability graphs}
\author{András Gyárfás, Márton Marits, Géza Tóth}
\date{arxiv:2408.13523v2; CC BY 4.0}
\begin{document}
\maketitle
\noindent\textit{Editorial note.} The counted unit is the article's own Theorem gensplit at source lines 225-226 together with the author's proof that opens with the bold header Proof. at source line 230 and closes with the square at source line 237; the author separates the two by an introductory sentence at source line 222 and blank lines, so the pairing is by adjacency through those lines only.\par\medskip
\noindent\textit{Editorial note.} The retained context at source lines 117-120 and 128-129 is the article's own description of comparability graphs and its definitions of p(G) and c(G); the retained context at source lines 204-206 and 208-209 is its own definition of a GSP graph and its own remark that GSP graphs are perfect and weakly triangulated. The sentences carrying the citations to Hayward and to the strong perfect graph theorem are deliberately not retained.\par\medskip
\noindent\textit{Editorial note.} The source preamble cannot be retained as it stands: it loads rawfonts, psfrag, verbatim, tikz with several libraries, hyperref with the pagebackref option and a dvi driver option for graphicx, none of which the retained body spans use, and it redefines backref commands that only exist because of that hyperref option. The retained preamble file therefore carries, line for line, source line 2, source line 3, source line 4, source line 16, source line 17, the theorem declarations of source lines 20-30 and the author macro definitions of source lines 32-41; the only textual change anywhere is the removal of the dvipdfm driver option from the graphicx line, because the assembly forces the xelatex engine. No new macro, package or declaration was introduced.\par\medskip
\noindent\textit{Editorial note.} No retained line contains a citation or a cross-reference, so the material carries no bibliography and needs no external navigation pointer.\par\medskip
\noindent\textit{Editorial note.} The article's Theorem twocover, Theorem promelsteger, Theorem bipline, Theorem unit, Theorem co-triangulated, Theorem loglog, Theorem ub, Lemma rec, the corollaries genupbound, shift, doblesh and p', and Proposition H9 are not counted here.\par\medskip
\noindent\textit{Allowed standard prerequisites (editorial).}\begin{itemize}
\item Comparability graph of a partially ordered set and the partition number p(G). A comparability graph has an edge between two vertices exactly when the corresponding elements of a partially ordered set are comparable, and p(G) is the minimum number of edge-disjoint comparability subgraphs of G whose union is the whole edge set; retained at source lines 117-120 and 128.. The counted theorem asserts p(G) <= 2, so both the class of graphs used in the partition and the parameter bounded must be available.
\item Generalized split graphs. A graph G is a GSP graph when either G or its complement can be written with its vertex set partitioned into V\_1 and V\_2, where V\_1 is a complete graph and V\_2 is a union of vertex-disjoint complete graphs with no edges between them; retained at source lines 204-206, and the article notes at source lines 208-209 that GSP graphs are perfect and form a subclass of the weakly triangulated graphs.. Supplies the two structural cases of the proof, one for G itself and one for its complement, and the hypothesis of perfectness used in the second case.
\end{itemize}\medskip
Here we focus on a well known family of perfect graphs, the {\em comparability graphs},
which can be obtained from a partially ordered set as follows.
Vertices correspond to the elements of the partially ordered set,
and two vertices are connected by an undirected edge if and only if the corresponding elements

Let $p(G)$ denote the minimum number $m$ such that $E(G)=\cup_{i=1}^m E(G_i)$ where the $G_i$-s are edge-disjoint comparability subgraphs of $G$.
Dropping  the condition that the $G_i$-s are edge-disjoint, we define $c(G)$ as the smallest number of comparability subgraphs of $G$ needed to {\em cover} $E(G)$.

\begin{definition} A graph $G$ is a GSP if either $G$ or its complement can be written in the following form:
  the vertex set is partitioned into $V_1,V_2$ where $V_1$ is a complete graph and $V_2$ is the union of vertex disjoint complete graphs with no edges between them.
\end{definition}

GSP graphs are perfect, form a subclass of weakly triangulated (alias weakly chordal) graphs,
defined by excluding induced cycles of length at least five and their complements.

\begin{theorem} \label{gensplit} For any GSP graph $G$, $p(G)\le 2$.
\end{theorem}

    {\bf Proof. } Assume first that the vertex set of a GSP graph $G$ is partitioned into $V_1,V_2$ where $V_1$
    is a complete graph and $V_2$ is the union of vertex disjoint complete graphs $K_i$ (with no edges between them).
    Then the complete subgraphs spanned by $V_1$ and by the $V(K_i)$-s form one comparability graph and the other is the bipartite graph
    $[V_1,V_2]$.

    In the second (complementary) case $V_1$ is an independent set in $G$ and $V_2$ spans a
    complete partite graph $H$, which is obviously a comparability graph.
    Then $H$ and the bipartite graph spanned by $[V_1,V(H)]$ in $G$ give the required partition.   \qed

\end{document}
